% !TEX root = ../main.tex
\section{Niklas}
Qualitative synthesis and theoretical framing rely heavily on accurate quotation and contextual integration.
Direct quotations allow researchers to anchor arguments in established literature while maintaining textual fidelity.

\subsection{Inline Quotations and Source Attribution}
In APA 7, short quotations (fewer than 40 words) are integrated directly into the text using double quotation marks.
The \texttt{\textbackslash enquote} command guarantees correct typographical curly quotes.
Every direct quotation must be accompanied by an exact page locator, paragraph number, or section header, such as \enquote{knowledge is built incrementally through repeated trial and iterative refactoring} \parencite[p.~48]{smith2023}.
When citing pageless online sources, cite the paragraph number instead, such as \parencite[para.~6]{smith2023}.

\subsection{Block Quotations and Qualitative Data}
Quotations containing 40 words or more must be presented in a freestanding block starting on a new line and indented 0.5 inches from the left margin.
In APA 7 block format, quotation marks are omitted, and the parenthetical citation is placed after the final closing punctuation:
\begin{quote}
Collaborative writing in version-controlled research environments presents unique workflow challenges.
When distributed teams establish semantic line breaking conventions and modular file divisions early in the project lifecycle, concurrent writing conflicts decrease dramatically and asynchronous review cycles become significantly more manageable. \parencite[p.~52]{smith2023}
\end{quote}
Following a block quotation, begin subsequent commentary on a new line with standard paragraph indentation if introducing a new topic, or flush left if continuing the same paragraph.
